\section{中断、设备与用户边界}

\subsection{从轮询到中断}

轮询由 CPU 主动反复读取状态；中断由设备在状态变化时请求处理。启用中断前，
内核必须完成 IDT、控制器屏蔽、EOI 规则和共享状态同步。中断上半部只确认和
记录事件，较长工作应延后执行。

\subsection{PIC、PIT、键盘与 IDE}

PIC 重映射避免 IRQ 与 CPU 异常向量冲突；PIT 提供周期时钟；键盘控制器提供
扫描码；IDE 提供扇区 I/O。每个驱动都应分为寄存器访问、状态机、策略和公共
接口，防止端口魔法数字散落在调度器或文件系统中。

v0.7 已完成这条最小链：关闭未配置虚拟线模式的本地 APIC，重映射双片
8259A，只开放 IRQ0/IRQ1；PIT 提供 64 位单调 tick，i8042 通过 IRQ1
交付集合 1 扫描码，内核 ATA PIO 重读 LBA 0。IRQ14、DMA、事件环形缓冲和
通用块层仍属于后续演进，不能因为“能读一个扇区”就提前宣称完成。

\subsection{Ring 3 与系统调用}

用户态不能直接访问内核页和特权指令。进入 Ring 3 需要用户代码/数据段、
用户栈、页表权限和正确的返回帧。系统调用入口必须先保存不可信上下文，再验证
编号、指针、长度和权限；验证通过才触碰内核对象。

\begin{contractbox}
用户指针不是 C++ 指针承诺，而是“可能映射、可能跨页、可能在验证后变化”的
不可信地址。复制进出用户空间必须定义部分失败和页故障处理。
\end{contractbox}

\subsection{驱动的完成语义}

命令写入设备不等于操作完成。驱动必须区分提交、设备接受、数据可用、完成确认
和错误恢复。这个模型会一直延伸到块层、文件系统和用户 \code{read} 返回值。

\subsection{中断替代高成本的持续查询}

轮询适合启动早期的单一设备，却让处理器在设备的时间尺度上空转。磁盘一次操作
可能相当于数百万条指令，键盘也只在人按键时产生数据。设备通过 IRQ 通知状态
变化后，处理器可以在等待期间运行其他任务，这成为多道程序与分时系统的硬件
基础之一。

中断并不是任意时刻直接调用普通函数。处理器先完成当前指令，在权限和屏蔽规则
允许时查找 IDT 门，保存最小返回现场，必要时切换特权栈，再跳到入口。汇编桩
保存 ABI 要求之外的状态，C++ 处理器确认中断来源并安排后续工作，最后恢复
现场。

\subsection{8259 PIC 的兼容角色}

传统双片 8259A 把 15 条可用 IRQ 汇聚到处理器。PC 初始向量与 CPU 保留异常
范围重叠，因此内核要把主片和从片重映射到独立向量区。初始化命令字的顺序、
级联线和屏蔽寄存器都属于设备状态机。

处理完成后必须发送 EOI，否则控制器会认为该优先级仍在服务，后续中断被阻塞。
来自从片的 IRQ 需要先确认从片，再确认主片级联。虚假 IRQ 则要按控制器在
服务寄存器判断，不能无条件访问不存在的设备。

后续 APIC 与 I/O APIC 提供更多向量、多处理器路由和更灵活触发方式，但项目
先在 PC 模型上完整实现 PIC，再把控制器差异封装在中断路由层。

\subsection{PIT 把时间变成可调度事件}

8254 PIT 使用固定输入时钟和可编程除数产生周期 IRQ。内核把每次 tick 累积为
单调时间，并在中断边界更新睡眠队列和调度预算。除数必须按 16 位硬件字段
写入，内部 tick 计数则使用 64 位，避免短时间回绕。

周期 tick 简单但会在空闲时持续唤醒 CPU。成熟系统会使用 APIC timer、
HPET 或 tickless 方案安排最近期限。本项目先用 PIT 建立时间与抢占，再在
抽象稳定后替换硬件来源。

\subsection{键盘扫描码不是字符}

PS/2 控制器返回按键按下与释放的扫描码序列。Shift、Ctrl 等修饰键改变解释
状态，扩展前缀还会组成多字节序列。驱动先把端口字节解析成规范按键事件，
键盘布局层再把按键与修饰状态转换为字符。把端口字节直接当 ASCII 会混淆
硬件位置、布局与文本编码。

IRQ 上半部只读取控制器数据并放入有界环形缓冲，避免设备缓冲溢出；终端层在
普通上下文消费事件。缓冲区满时要有明确丢弃或背压策略，不能静默覆盖仍未处理
的数据。

\subsection{中断上下文的并发限制}

中断可能打断持锁代码。若处理函数再次获取同一把锁，就会在单核上永久等待自己。
锁设计因此要说明中断上下文是否使用、获取时是否关闭本地中断、能否睡眠以及
全局锁顺序。共享数据发布还要有编译器和处理器内存顺序。

上半部完成确认设备、抓取最小数据和唤醒工作；耗时解析、复制与文件系统操作
延后到可调度上下文。这个拆分同时缩短中断关闭时间和最坏响应延迟。

\subsection{特权级建立用户程序的硬边界}

x86 用 Ring 0 到 Ring 3 表示特权，常规操作系统使用 Ring 0 运行内核，
Ring 3 运行应用。页表 U/S 位限制用户访问内核页，指令权限限制
\code{in}、\code{out}、控制寄存器和中断控制。仅把代码标签称为“用户程序”
并不会产生隔离，必须同时建立代码段、用户栈、页权限和受控入口。

首次进入用户态可构造 \code{iretq} 帧，包含用户 SS、RSP、RFLAGS、CS 和
RIP。此后系统调用或异常使处理器切回内核入口，TSS 提供可信内核栈。内核不能
继续使用用户栈处理特权数据。

\subsection{系统调用把不可信请求转换成内核操作}

系统调用号选择操作，寄存器携带标量或用户地址。入口先保存用户上下文，切换
到内核数据模型，再验证编号、长度、算术溢出和页权限。跨页缓冲可能前半合法、
后半无效，所以复制函数必须定义部分完成语义。

验证后仍不能长期保存裸用户指针，因为另一个线程可能改变映射。内核通常在受控
窗口内复制数据，或固定页面并持有地址空间锁。所有返回值使用稳定 ABI，不能
把内部 C++ 对象布局暴露给用户程序。
